%% ============================================================
%% 分片 32：原文第 11 节「$A$-Branes and ${\cal D}$-Modules」（原文 glshort.tex 起于第 9938 行）
%% 由 tools/harvmac2tex.py 从 arXiv 源码（hep-th/0604151v3）自动生成。
%% 只译英文散文；公式、编号、\label、\cite 一律不改。
%% ============================================================

\subsection{典范余各向同性 $A$-膜}

我们的出发点是 \cite{ref80} 中的一个观察：对一般的辛流形 $X$，可能存在
不支撑在 Lagrange 子流形上、而支撑在更一般的余各向同性子流形上的
$A$-膜。辛流形 $X$ 的余各向同性子流形 $Y$ 是指局部由第一类约束定义的
子流形；换言之，它局部由若干 Poisson 对易的函数 $f_i$，$i=1,\dots,d$
的零点集定义。$X$ 上独立的 Poisson 对易函数的最大数目是 $X$ 维数的一半，
所以 $Y$ 的维数至少是 $X$ 维数的一半。若 $Y$ 的维数恰好是 $X$ 维数的
一半，则 $Y$ 是 Lagrange 的，而支撑在 $Y$ 上的 $A$-膜就是通常的
Lagrange $A$-膜。支撑维数高于中维的余各向同性 $A$-膜一般也存在，不过
它们存在的条件相当特殊。

我们只需要 $Y=X$ 且 $A$-膜秩为一的情形。于是其 Chan-Paton 丛是
$U(1)$ 丛 ${\cal L}\to X$。Chan-Paton 规范场取值于 $U(1)$ 的李代数，
而这是一个一维实向量空间。因此该规范场局部就是普通的实一形式 $A$，
其协变微分为 $D=d+A$；曲率定义为 $F=dA$。在 $U(1)$ 的酉表示中，
李代数 $\frak{u}(1)$ 的元素 $1$ 通过 $1\times 1$ 的反厄米矩阵 $-in$
作用，$n$ 为某个整数，而平行移动为 $\exp(in\int A)$。我们把 $X$ 上的
Chan-Paton 丛解释为与表示 $n=1$ 相伴的线丛。对 sigma 模型的映射
$\Phi:\Sigma\to X$，世界面路径积分中与 Chan-Paton 丛相关的因子，是
$\Phi^*(A)$ 沿 $\partial\Sigma$ 在 $n=1$ 表示中的和乐（holonomy），即
$\exp(i\oint_{\partial\Sigma}\Phi^*(A))$。

对一个 $Y=X$、Chan-Paton 线丛为 ${\cal L}$ 的膜，要成为 $A$-膜，其曲率
$F$ 必须满足
% 式 (11.1)
\setcounter{equation}{0}
\begin{equation}\label{eq:coisobrane}
(\omega^{-1}F)^2=-1.
\end{equation}
 也就是说，$(1,1)$ 张量 $\omega^{-1}F$ 是一个近复结构。形式 $F$ 与
$\omega$ 都是闭的，与 \eqref{eq:coisobrane} 一起可推出近复结构
$\hat N=\omega^{-1}F$ 是可积的 \cite{ref80}。容易看出，$\omega$ 与 $F$
关于 $\hat N$ 都是 $(2,0)\oplus(0,2)$ 型的。此外，复形式 $\omega-iF$
是 $(2,0)$ 型的，因而是 $X$ 上的一个全纯辛形式。所以这类余各向同性
$A$-膜只有在 $\omega$ 是某个全纯辛形式的实部时才可能存在。关于在
广义复几何中构造这类膜，另见 \cite{ref67}。

我们感兴趣的是 $\MH(G,C)$ 上辛结构为 $\omega=(\Im\tau)\,\omega_K$ 的
$A$-模型。此时式 \eqref{eq:coisobrane} 有一个显然的解： % 式 (11.2)
\setcounter{equation}{1}
\begin{equation}\label{eq:yetro}
F=
(\Im\tau)\,\omega_J,\quad \hat N=\omega_K^{-1}\omega_J= I.
\end{equation}

（也存在把 $F$ 与 $\hat N$ 的符号反过来的第二个解。）二形式 $\omega_J$
是恰当的，正如我们在第~\ref{sec:mhhyper}~节中所解释的，所以 $F$ 是平凡
线丛 ${\cal N}\to\MH(G,C)$ 上某个联络的曲率二形式。联络形式为 % 式 (11.3)
\setcounter{equation}{2}
\begin{equation}\label{eq:exactcc}
\alpha={\Im\tau\over {2\pi}}\int_C |d^2z|\,
\Tr \left(\phi_z\delta A_\bz +\phi_\bz \delta A_z\right)
\end{equation}
 且 $F=d\alpha$。虽然更一般的二形式 % 式 (11.4)
\setcounter{equation}{3}
\begin{equation}\label{eq:lombo}
(\Im\tau)\left(a\omega_J+b\omega_I\right), \quad a^2+b^2=1,
\end{equation}
 也满足式 \eqref{eq:coisobrane}，但它在 $b\neq 0$ 时不是恰当的，
并且对一般的 $\tau$ 值，其周期不是 $2\pi$ 的整数倍。所以在一般情况下，
这个形式不是任何联络的曲率二形式，不能用来构造 $A$-膜。

我们把 $\MH$ 上满足 $F=(\Im\tau)\,\omega_J$ 的 $A$-膜称为典范余各向同性
膜，简称 c.c. 膜。我们有时也把它记作 ${\cal B}_{c.c.}$。若 $G$ 单连通，
则 $\MH(G,C)$ 也单连通，此时 c.c. 膜是唯一的；否则，它唯一到把其
Chan-Paton 丛扭化一个平坦线丛为止。

关于 c.c. 膜的第一个观察是：它是一个 $(A,B,A)$-膜，即在复结构 $I$ 与
$K$ 中是 $A$-膜、在复结构 $J$ 中是 $B$-膜。要验证它在复结构 $I$ 中是
$A$-膜，我们只需要恒等式
% [WARN] 公式 \eqn\eqref{eq:hombo} 无编号或重复
\begin{equation*}
\tag*{(11.5)}
(\omega_I^{-1}\omega_J)^2=-1.
\end{equation*}
 这意味着 c.c. 膜在复结构 $I$ 中满足余各向同性膜的条件，正如它在复结构
$K$ 中一样。要验证它在复结构 $J$ 中是 $B$-膜，只需注意膜上规范场的曲率
是 $\omega_J$ 的倍数，而 $\omega_J$ 在复结构 $J$ 中是 $(1,1)$ 型的。

这些公式其实只在 $\MH$ 的光滑部分定义了 c.c. 膜。不过，c.c. 膜可以由
底层规范理论中的边界条件来定义；详见第~\ref{sec:xgoto}~节。这一论断在
量子场论上等价于说：该构造在整个 $G$-丛的「栈」上都成立。

接下来，我们想理解两端都落在 c.c. 膜上的拓扑开弦。它们由 BRST 不变的
边界算子插入来代表。一般而言，正如我们在第~\ref{sec:comptwo}~节 的
\dunamp\ 中所概述的，对任何膜 ${\cal B}$，$({\cal B},{\cal B})$ 弦构成
一个结合代数（但未必交换）。用数学语言说，这就是把膜 ${\cal B}$ 视为
$\MH(G,C)$ 上膜范畴的一个对象时的自同态代数。我们想对 c.c. 膜的情形
计算这个自同态代数。

作为第一步，我们要论证世界面瞬子对计算没有贡献。在 sigma 模型中定义
余各向同性膜的边界条件 \cite{ref80} 说：端点落在像我们这里所考虑的、
填满整个空间的余各向同性膜上的弦，必须满足 % 式 (11.6)
\setcounter{equation}{5}
\begin{equation}\label{eq:melbo}
G_{IJ}\partial_1\Phi^J=-i
F_{IJ}\partial_0\Phi^J,
\end{equation}
 其中 $\partial_1$ 与 $\partial_0$ 分别是世界面边界处的法向导数与切向
导数。再施加瞬子方程（即映射 $\Phi$ 应在复结构 $K$ 中全纯），以便把
$\partial_1\Phi$ 与 $\partial_0 \Phi$ 联系起来，我们得到 % 式 (11.7)
\setcounter{equation}{6}
\begin{equation}\label{eq:zimbolo}
(\omega-iF)\partial_0\Phi=0.
\end{equation}
 由于 $\omega$ 与 $F$ 都是非退化的，瞬子在边界上必须是常数。但由
解析性，它必须在处处为常数，从而是平凡的。

既然没有瞬子修正，有趣的东西就都必须来自世界面微扰论。但世界面微扰论
在 $\hbar$ 的任意阶都会生成靶空间局部的 sigma 模型算子与耦合的修正。
因此，至少在世界面微扰论的所有阶，很自然要考虑只在靶空间中局部定义的
边界可观测量。这样的可观测量构成靶空间上的一个层，于是我们得到的不是
可观测量代数，而是在微扰论所有阶都更强的东西：一个代数层。（在闭弦
情形，这种靶空间中的局部化曾在 \cite{ref133,ref134,ref135} 中讨论过，
那里解释的是更早的数学工作 \cite{ref136}。）在微扰论中分析这个层的
同时，我们也会讨论精确理论中会发生什么。

对任意余各向同性 $A$-膜 $Y$，BRST 不变的边界算子空间已在经典极限下于
\cite{ref80,ref132} 中确定。就我们的目的而言，只需考虑鬼数为零的边界
算子。于是在 $Y=X$ 的情形，这样的可观测量就是 $X$ 上的全纯函数，其中
$X$ 被看作带有复结构 $\hat N$ 的复流形。特殊化到这一情形并在靶空间中
局部化，我们得到：作为加性群（忽略环结构），c.c. 膜的拓扑边界算子层
就是 $\MH(G,C)$ 上关于复结构 $I$ 全纯的函数层。注意，虽然我们在复结构
$K$ 的 $A$-模型中工作，但这里相关的复结构是 $\hat N=I$。局部地，边界
算子层没有更高阶上同调，也不能形变，因此不可能有量子修正。但整体上它
有量子修正，我们将对此进行分析。

 \cite{ref137}  $\MH$ 上的全纯函数层有一个显然的交换代数结构，这正是
经典极限下正确的结构。在「物理的」弦理论背景下，当规范场强 $F$ 存在
时，开弦代数会变得非交换 \cite{ref137}。

让我们先非正式地讨论一下这在当前语境下会如何运作。通常的非交换参数
（在 \cite{ref137} 中称为 $\theta$）是 $(g+F)^{-1}$ 的反对称部分。在
当前语境中，它成为 $\omega_J^{-1}$ 的倍数。$\omega_J$ 是 $\Omega_I$ 的
实部，而 $\Omega_I$ 是复结构 $I$ 中的全纯 $(2,0)$-形式。当我们相乘由
关于复结构 $I$ 全纯的函数（它们正是拓扑可观测量）所导出的可观测量时，
$\omega_J^{-1}$ 可以换成 $\Omega_I^{-1}$。所以这类可观测量的乘法按照
由 $\Omega_I$ 导出的 Poisson 括号来形变，这一点我们马上就会看到。

在物理弦理论中，对任何 $F\not=0$，开弦的乘法都含有来自 $(g+F)^{-1}$ 的
非交换相位，但一般还有许多其他弦论效应。用靶空间上函数代数的非交换
形变来精确描述，通常只在 $|F|>>|g|$ 时才出现 。在当前
问题中我们并不处于这个极限；相反，我们有 $|F|=|g|$。不过，一旦限制到
拓扑可观测量，我们就会得到一种简化，类似于物理弦理论中 $|F|>>|g|$ 时
通常发生的情形。在 $B$-型拓扑弦的语境中，这一点最早在 \cite{ref71} 中
讨论过。实质上，物理弦理论中只有在 $|F|>>|g|$ 时才发生的弦激发消除，
在这里是由于限制到拓扑可观测量而出现的。

现在我们转向拓扑理论的更精确分析。圆盘上 $A$-模型的作用量可以写成
如下形式 % 式 (11.8)
\setcounter{equation}{7}
\begin{equation}\label{eq:mortop}
S=\int_\Sigma \Phi^*(\omega-iF)+\int_\Sigma d^2\sigma \{Q,V\},
\end{equation}
 其中 $Q$ 是 BRST 算子或拓扑超荷。在我们的情形中，形式
$\omega-iF=-i \Omega_I \Im\tau$ 是恰当的（当然这一点局部总是成立）。
若引入 $\Omega_I$ 的势 $\varpi$，使得 $\Omega_I=d\varpi$，则作用量取
形式 % 式 (11.9)
\setcounter{equation}{8}
\begin{equation}\label{eq:ortop}
S=-i\,\Im\tau\int_{\partial\Sigma}\Phi^*(\varpi) +\int_\Sigma
d^2\sigma \{Q,V\}.
\end{equation}
 此外，$\varpi$ 是 $\M_H(G,C)$ 上典范全纯一形式的倍数： % 式 (11.10)
\setcounter{equation}{9}
\begin{equation}\label{eq:honop}
\varpi={1\over \pi}\int |d^2 z|\, \Tr \left(\phi_z\delta
A_\bz\right).
\end{equation}
  若限制到 $T^*\M(G,C)$，我们可以写得更具体。若 $q^\alpha$ 是
$\M(G,C)$ 上的局部全纯坐标，$p_\alpha$ 是 $T^*\M(G,C)$ 的纤维上与
$\phi_\alpha$ 典范共轭的线性函数，则
$\varpi=\sum_\alpha \,p_\alpha\,dq^\alpha$，且
$S=-i\,{\rm Im}\,\tau\int p_\alpha\,dq^\alpha+\{Q,\dots\}$。于是，
在模去 $Q$-恰当项的意义下，$A$-模型的作用量就是粒子在 $\cM(G,C)$ 上
运动的量子力学作用量的全纯版本（哈密顿量为零）。

我们无疑会因特殊化到 $T^*\M\subset \MH$ 而失去一些重要结构，因为
c.c. 膜是在整个 $\MH$ 上定义的，甚至在四维规范理论中也不只是在
$T^*\M$ 上定义。但让我们看看作了这一特殊化之后还剩下什么结构。

式 \eqref{eq:ortop} 中作用量的形式是标准的，结果也是熟悉的：在
${1/ {\rm Im}\,\tau}$ 的一阶，微扰的效果是按照由该作用量导出的
Poisson 括号来形变局部定义的全纯函数的交换代数。所以在
${1/{\rm Im}\,\tau}$ 的微扰论中，c.c. 弦（更严格地说，
$({\cal B}_{c.c.},{\cal B}_{c.c.})$ 弦）构成非交换代数的层。但究竟
得到哪个层？

一个有力的工具是本文第~\ref{sec:complexstr}~节 末尾引入、并在后文多处
使用的 $\CC^*$ 群 ${\cal U}$。它按 % 式 (11.11)
\setcounter{equation}{10}
\begin{equation}\label{eq:utyp}
\phi_z\to \lambda
\phi_z
\end{equation}
 作用，而可观测量空间上相应的分次曾在第~\ref{sec:sdualhitfib}~节 中
用来分析 $S$-对偶对零膜的作用。$\varpi$ 在这一标度变换下的次数为 1；
因此，若把次数 $1$ 赋给 ${1/{\rm Im}\,\tau}$，也就是赋给相应量子力学
问题中的普朗克常数 $\hbar$，就可以把 ${\cal U}$ 提升为量子理论的
对称性。

在经典层面，可观测量的乘法保持 ${\cal U}$ 的分次；次数为 $m$ 与 $n$
的可观测量之积是次数为 $m+n$ 的可观测量。在量子层面，乘积中还可能
出现次数更低的项，只要它们带上普朗克常数的正幂。因此量子可观测量的
代数只是带滤过的，而不是分次的。

为了有效地利用 $\Bbb{C}^*$ 对称性来分析这个代数，我们需要考虑在
$\Bbb{C}^*$-不变开集上定义的可观测量。为此，取开集
${{\cal V}}\subset {\EUM}$，并考虑定义在 $T^*{{\cal V}}\subset \MH$ 上
的可观测量。于是我们只对端点落在 c.c. 膜上的弦的空间作部分层化，只
允许 $\MH$ 中这种形式的开集。这种部分层化的好处是它非微扰地成立，
这一点我们马上就会看到。若对任意开集 ${\cal W}\subset \MH$ 都定义
可观测量的结合代数，这样完全层化在 $1/{\rm Im}\,\tau$ 的所有阶都运行
良好，但指望它在精确理论中也成立是不现实的。关于这类问题的普遍讨论，
见 \cite{ref138}。

为了理解为什么一般存在阻碍，使形变量子化不能超出微扰论来层化，考虑
一个简单例子：坐标为 $x,y$、辛形式为 $dx\wedge dy$ 的 $\Bbb{C}^{2}$ 上
的形变量子化。此时若把自己限制在 $\Bbb{C}^{2}$ 上的多项式函数，形变
量子化可以用 Moyal-Wigner 乘积精确定义\footnote{他们的论文重印本、许多
其他原始论文以及一篇综述，见 \cite{ref139}。}。要层化这一情形，需取
开集 ${\cal W}\subset\Bbb{C}^{2}$，并试图定义在 ${\cal W}$ 中全纯的
函数的 Moyal-Wigner 乘积。到 $\hbar$ 的任何有限阶都没有问题，但
Moyal-Wigner 乘积涉及无穷求和这一事实，使得要得到精确结果就出了问题。
为看出问题所在，尝试在原点计算函数 $1/(x-a)$ 与 $1/(y-b)$ 的
Moyal-Wigner 乘积，二者各自在 $\Bbb{C}^2$ 中适当的开集上全纯。

所以我们把自己限制在部分层化：考虑 $\MH$ 中形如 $T^*{{\cal V}}$ 的开集，
其中 ${{\cal V}}\subset\M$。此外，我们只考虑 $T^*{{\cal V}}$ 上沿投影
$T^*{{\cal V}}\to{{\cal V}}$ 的纤维有多项式增长的函数。这样的函数构成
的代数由次数小于或等于 $1$ 的元素生成，也就是由沿纤维至多线性的函数
生成。

次数为 $0$ 的元素就是 ${{\cal V}}$ 上的全纯函数。这类函数的乘法代数
是无法形变的，因为 $T^*{{\cal V}}$ 中没有负次数的可观测量。由于
$\hbar$ 的次数为 1，对次数 $0$ 的函数的乘法作修正，必须涉及一个
负次数的可观测量。

次数为 $1$ 的元素是余切丛纤维上的线性函数。于是若 $q^\alpha$ 是
${{\cal V}}$ 上的局部坐标，$p_\alpha$ 是纤维上与之共轭的线性函数，
则次数 $1$ 的元素形如 $\kappa=\sum_\alpha p_\alpha\, V^\alpha(q)$，
并与 $\cM$ 上的全纯向量场 $\sum_\alpha
V^\alpha{\partial\over\partial q^\alpha}$ 自然对应。在相乘一个次数
$0$ 的对象 $f$ 与一个次数 $1$ 的对象 $\kappa$ 时，只有「空间」容纳
一阶形变——因为乘积中 $\hbar^2$ 阶的项又会乘上负次数的算子。一阶形变
由 Poisson 括号刻画 % 式 (11.12)
\setcounter{equation}{11}
\begin{equation}\label{eq:hembo}
[\kappa,
f]=-i(\Im\tau)^{-1}\kappa(f).
\end{equation}

为继续下去，我们用 $\hat w$ 表示对应于 $T^*{{\cal V}}$ 上全纯函数 $w$
的开弦态。公式 \eqref{eq:hembo} 决定了对应于 $f$ 与 $\kappa$ 的环元素
$\hat f$、$\hat\kappa$ 的 $\star$ 乘积（即开弦乘法的代数），但并不完全
唯一，因为（在保持代数滤过方式不变的前提下）我们可以给 $\hat\kappa$
加上一个次数为零的可观测量。若对 $\kappa=V^\alpha(q) p_\alpha$，定义
$\hat\kappa =\hat V^\alpha\star \hat p_\alpha$，则代数就唯一确定了，它
就是 % 式 (11.13)
\setcounter{equation}{12}
\begin{equation}\label{eq:yembo}
\begin{aligned}\hat f\star\hat\kappa&
=\hat{f\kappa}\\ \hat\kappa\star \hat f&
 =\hat{f\kappa}-i(\Im\tau)^{-1}\hat{\kappa(f)}.\\\end{aligned}
\end{equation}

 当然，$\hat{ f\kappa}$ 与 $\hat{\kappa(f)}$ 是对应于函数 $f\kappa$ 与
$\kappa(f)=V^\alpha\,\partial f/\partial q^\alpha$ 的环元素。我们还没有
完全用完重新定义 $\hat\kappa$ 的自由；剩下的自由是：对余切丛上的线性
函数 $p_\alpha$，我们可以作替换 % 式 (11.14)
\setcounter{equation}{13}
\begin{equation}\label{eq:remob}
\hat p_\alpha\to \hat p_\alpha+\hat A_\alpha
\end{equation}
 其中 $A_\alpha(q)$ 是某些函数。

当我们相乘两个次数均为 $1$ 的函数 $\kappa$ 与 $\kappa'$ 时，会出现
新东西。有一阶形变，它仍由 Poisson 括号给出，但现在也有容纳二阶形变
的空间。$p_\alpha$ 与 $p_\beta$ 的 Poisson 括号为零，但 $\star$ 代数
中的对易子 $[\hat p_\alpha,\hat p_\beta]$ 可能等于
$\hbar^2 F_{\alpha\beta}(q)$，其中 $F_{\alpha\beta}$ 是某些函数。利用
代数的结合性可以证明，二形式
$F_{\alpha\beta}dq^\alpha\,dq^\beta$ 是闭的，因而（对小开集
${{\cal V}}$）可以通过变换 \eqref{eq:remob} 置为零。

因此，对小开集 ${\cal V}\subset \MH$，与 $T^*{{\cal V}}$ 相伴的代数
同构于 ${{\cal V}}$ 上的全纯微分算子代数，其中 $\M$ 上的函数按乘法
作用，且 % 式 (11.15)
\setcounter{equation}{14}
\begin{equation}\label{eq:mebolo}
\hat p_\alpha=-{i\over {\rm
Im}\,\tau}{\partial\over\partial q^\alpha}.
\end{equation}

现在假设用开集 ${{\cal V}}^{(i)}$ 覆盖 $\M$。在重叠区域
${{\cal V}}^{(i)}\cap {{\cal V}}^{(j)}$ 中，次数为零的环元素 $\hat f$
必须一致，但次数 $1$ 的元素 $\hat p_\alpha^{(i)}$ 与
$\hat p_\alpha^{(j)}$ 可以相差一个变换 \eqref{eq:remob}。由于我们在
每个开集中要求 $[\hat p_\alpha,\hat p_\beta]=0$，它们必须通过
% 式 (11.16)
\setcounter{equation}{15}
\begin{equation}\label{eq:hecox}
\hat p_\alpha^{(i)}=\hat
p_\alpha^{(j)}+{1\over {\rm Im}\,\tau}\hat{\partial s^{(ij)}\over
\partial q^\alpha}
\end{equation}
 相联系，其中 $s^{(ij)}$ 是 ${{\cal V}}^{(i)}\cap {{\cal V}}^{(j)}$ 上
的某些局部全纯函数。这一关系的相容性推出：在三重重叠中，量
$c^{(ijk)}=s^{(ij)}+s^{(jk)}+s^{(ki)}$ 是复常数。此外，它们自动满足
上闭链条件，并定义了 $H^2({\EUM},\Bbb{C})$ 的一个元素。

若这些常数取值于 $2\pi\Bbb{Z}$，则量 $\exp(-is^{(ij)})$ 是复线丛
${\cal L}\to \MH$ 的转移函数。此时，开弦代数层就是作用在 ${\cal L}$ 的
截面上的微分算子层。我们把这个代数层记作 ${\cal D}_{\cal L}$。

${\cal L}$ 并不完全唯一。一般地，代数层 ${\cal D}_{\cal L}$ 在把
${\cal L}$ 扭化一个平坦线丛之下不变，因为平坦线丛有常数转移函数，
与动量 $\hat p_\alpha$ 对易。

没有理由认为常数 $c^{(ijk)}$ 必须适当量子化，即必须取值于
$2\pi\Bbb{Z}$。例如，从一个导致线丛 ${\cal L}$ 上微分算子层
${\cal D}_{\cal L}$ 的适当量子化的情形出发，我们可以把 $s^{(ij)}$ 乘以
一个复数 $\gamma$，从而定义一个新的层 ${\cal D}_{{\cal L}^\gamma}$。
若 ${\cal L}$ 拓扑非平凡，则线丛 ${\cal L}^\gamma$ 不存在（$\gamma$ 取
整数或也许有理数时除外），但相应的代数层
${\cal D}_{{\cal L}^\gamma}$ 确实存在。

由于 $H^2({\EUM}(G,C),\Bbb{C})$ 是一维的，在这个具体例子中，由这类
构造可能得到的最一般的代数层形如 ${\cal D}_{{{\frak L}}^\gamma}$，
其中 $\gamma$ 是复数，${{\frak L}}$ 是 $\EUM$ 上的基本线丛——我们
笼统地称之为行列式线丛。如果把 $\EUM$ 换成更一般的空间 $X$，则还可能
有略微的推广，涉及线丛复幂的乘积。

无论如何，正如我们将看到的，几何 Langlands 纲领的基本情形涉及某个
普通线丛上的微分算子。但我们在第~\ref{sec:genotics}~节 中考虑的一个
推广确实涉及线丛的复幂。

注意，在上述分析中，不存在微扰论是否收敛的问题，因为一切都由二阶发生
的事情决定。倘若我们试图完全层化，即对 $\MH$ 中每个开集都配上一个
代数，我们就会在微扰论的收敛性上遇到麻烦。

\par\medskip\emph{时间反演对称性}

这些多少有些抽象的想法并不能完全确定 c.c. 弦的代数层。它是某个线丛
（或线丛的复幂）上的微分算子层，但究竟是哪一个？为确定它，我们将利用
时间反演对称性。从时间反演对称性得到额外约束是很自然的，因为正如我们
在式 \eqref{eq:zongor} 中看到的，它恰好在我们目前讨论所用的值 $t=1$
处才是扭化拓扑场论的对称性。

在规范理论中，时间反演（我们记作 ${\EUT}$）反转四维流形 $M$ 的时间
坐标，而在空间坐标上平凡地作用。若 $M=\Bbb{R}\times I\times C$，则
${\EUT}$ 就通过 $\Bbb{R}$ 上的 $-1$ 作用。在有效二维理论中，${\EUT}$
是世界面 $\Sigma$ 的一个定向反转对称，它把 $\Sigma$ 的边界映到自身，
同时反转边界的定向。由于我们的考虑是局部的，可以假定 $\Sigma$ 或 $M$
具有这样的对称性。

在 ${\EUT}$ 下，玻色子 $(A,\phi)$ 变换为
$({\EUT}^*A,-{\EUT}^*\phi)$，正如我们在式 \eqref{eq:yoro} 中看到的。
若 ${\EUT}$ 把 $\phi$ 映为 ${\EUT}^*\phi$，则 \eqref{eq:mortop} 与
\eqref{eq:ortop} 中涉及 $\omega$、$F$ 与 $\varpi$ 的项在世界面定向
反转下会是奇的。但有了额外的负号，这些项就不变了。

时间反演反转世界面边界上算子插入的顺序，因此它把 c.c. 弦的代数映到其
「相反」代数，其中元素相同但乘法顺序相反。于是对代数 ${\eurm C}$ 的
每个元素 $x$，相反代数 ${\eurm C}^{op}$ 有元素 $x^{op}$，满足
$(x\star y)^{op}=y^{op}\star^{op} x^{op}$。这里 $\star$ 是
${\eurm C}$ 中的乘法，$\star^{op}$ 是 ${\eurm C}^{op}$ 中的乘法。

复流形 $X$ 上线丛 $\cL$ 上的微分算子代数的相反代数，同构于
$\cL^{-1}\otimes K_X$ 上的微分算子代数 \cite{ref140}，其中 $K_X$ 是
$X$ 的典范线丛（即最高次全纯形式丛）。为看出这一点，令 $s$ 与 $u$
分别为 $\cL$ 与 $\cL^{-1}\otimes K_X$ 的紧支撑截面（不必全纯），并令
${\eurm F}$ 是作用在 $\cL$ 上的全纯微分算子。则
$s\, {{\eurm F}}u$ 是 $K_X$ 的截面。$K_X$ 的截面是中间维数的微分形式，
可在 $X$ 的任何实切片 $Z$ 上积分。此外，通过分部积分，可以构造
「转置」算子 ${{\eurm F}}^t$，它作用在 $\cL^{-1}\otimes K_X$ 的截面上，
并对所有 $u$ 与 $s$ 满足
$\int_Z u\,{{\eurm F}} s=\int_Z ({\eurm F}^t u)s$。${{\eurm F}}^t$ 与
${{\eurm F}}$ 一样是带全纯系数的微分算子，且不依赖于 $Z$ 的选取。
映射 ${{\eurm F}}\to {{\eurm F}}^t$ 反转乘法顺序，这证明
${\cal D}_{\cal L}$ 的相反代数层是
${\cal D}_{{\cal L}^{-1}\otimes K_X}$。

c.c. 膜的时间反演对称性意味着：由端点落在 c.c. 膜上的开弦得到的代数层
必须同构于其相反代数。这一点，再加上它是作用在某个线丛上的微分算子层
这一事实，推出该代数必须是 ${\cal D}_{K_{\EUM}^{1/2}}$，即作用在
$K_{\EUM}^{1/2}$（$\M$ 的典范丛的平方根）上的微分算子层。$\M$ 实际上
是自旋流形，所以 $K_\M^{1/2}$ 是普通线丛。$\M$ 的自旋结构未必唯一，
但由于作用在线丛上的微分算子层在扭化一个平坦线丛之下不变，代数层
${\cal D}_{K_\EUM^{1/2}}$ 不依赖于 $K_{\EUM}^{1/2}$ 的选取。

\par\medskip\emph{拓扑弦的整体代数}

现在我们要讨论\emph{整体}的 c.c. 弦。

如何从可观测量层得到整体可观测量的代数？只需取可观测量层的整体截面，
或者更一般地，取 $\M$ 的取值于该层的上同调。为简单起见，我们只考虑
整体截面，即鬼数或上同调次数为零（但对 $\Bbb{C}^*$ 对称 ${\cal U}$ 有
任意标度次数）的可观测量。

在经典层面，可观测量层就是 $\MH$ 上的全纯函数层，其整体截面就是
$\MH$ 上的整体全纯函数。它们就是可积系统的对易哈密顿量，我们在
第~\ref{sec:hitchfib}~节中描述过。

同样地，在量子层面，要得到整体可观测量，我们必须取局部可观测量层的
整体截面。但定义这个层的粘接律相对于经典情形有量子修正。这一量子修正
在我们的分析中明确地出现在式 \eqref{eq:hecox} 中。所以对整体可观测量的
描述可能有量子修正。

既然已经理解了 c.c. 弦的层是 $\M$ 上的微分算子层，我们就可以具体地
理解这一修正意味着什么。设给定 $\MH$ 上的一个全纯函数，例如一个次数
为三的函数，一旦在 $\M$ 上取定局部复坐标，它看起来像
$w=\sum_{\alpha,\beta,\gamma} f^{\alpha\beta\gamma}(q)p_\alpha
p_\beta p_\gamma$。要把它提升为整体的 $({\cal
B}_{c.c.},{\cal B}_{c.c.})$ 弦，我们需要找到 $\M$ 上一个三阶微分算子
$\hat w$，其「符号」为 $w$。换句话说，我们希望 $\hat w=(-i/{\rm
Im}\,\tau)^3\,f^{\alpha\beta\gamma}\partial^3 /\partial
q^\alpha\partial q^\beta\partial q^\gamma+\dots$，其中省略号指二阶或
更低阶的微分算子。由于 $\M$ 上没有典范的局部复坐标 $q^\alpha$，这样
的算子局部并不唯一，整体上也可能不存在。

若符号为 $w$ 的全纯微分算子 $\hat w$ 确实存在，我们就说函数 $w$ 可以
「量子化」。$\MH$ 上的对易哈密顿量都可以量子化。对次数为二的函数，这
最先由 Hitchin \cite{ref141} 证明；对任意次数的函数，则由 Beilinson 与
Drinfeld \cite{ref5} 证明；这一结果也是他们论文标题的理据之一。\footnote{他们的论证基于二维共形场论，更确切地说，基于群 $G$
在水平 $k=-h$ 的流代数，其中 $h$ 是 $G$ 的对偶 Coxeter 数。} 我们将
给出一个论证，其中进一步用到时间反演对称性 ${\EUT}$。这一论证还表明
所得的微分算子相互对易，这一点最先在上述文献中证明。

我们用到的事实是：$({\cal B}_{c.c},{\cal B}_{c.c.})$ 弦对应于可以
插入到由 c.c. 膜标记的边界上的边界算子。此外，这些边界可观测量可以
从四维规范理论的算子导出，这是一个有力的工具，因为它具有局部性，而
这种局部性在约化到二维时会丢失。这将使我们能给出一个简单的论证。
事实上，正如我们在第~\ref{sec:hitchfib}~节中回顾的，对易哈密顿量形如
$\int_C \alpha\,{\cal P}(\varphi)$，其中 ${\cal P}$ 是某个次数 $k$ 的
齐次规范不变多项式，$\alpha\in H^1(C,K_C^{1-k})$。要证明这样的算子
$Q$-不变，只需证明局部算子 ${\cal P}(\varphi)$ 是 $Q$-不变的。我们在
第~\ref{sec:xgoto}~节中通过分类可能的算子并使用时间反演对称性来做这
件事。从这个观点看，显然通过量子化两个经典表达式
$\int_C \alpha \,{\cal P}(\varphi)$ 与 $\int_C\tilde
\alpha \,\tilde {\cal P}(\varphi)$ 得到的微分算子相互对易，因为可以取
$\alpha$ 与 $\tilde\alpha$ 有不相交的支集。

为了说明「$\MH$ 上的经典全纯函数可以量子化为 $K_\M^{1/2}$ 上的微分
算子」这一事实的非平凡性，让我们指出 \cite{ref141,ref5}：它们不能
量子化为除 $K_\cM^{1/2}$ 以外任何线丛上的微分算子。它们能量子化为
\emph{相互对易}的微分算子这一事实意味着：$\MH$ 上 Hitchin 的经典
可积系统可以「量子化」为 $\M$ 上的量子可积系统。

\begin{yizhu}[易错点]
按：式 \eqref{eq:coisobrane} 与式 \eqref{eq:yetro} 是同一件事的两种写法，
符号容易看错。在复结构 $K$ 中，$\hat N=\omega_K^{-1}\omega_J=I$ 本身就是
近复结构，故 $(\omega^{-1}F)^2=-1$ 自动成立（此即 (11.2)；注意这里的
$I$ 是复结构，$I^2=-1$，不要读成恒等矩阵）。在复结构 $I$ 中作 $A$-模型
时，辛形式取 $(\Im\tau)\omega_I$，而 $F=(\Im\tau)\omega_J$，于是
$\omega^{-1}F=\omega_I^{-1}\omega_J$，$\Im\tau$ 因子相消，余各向同性条件
正是 (11.5) 的 $(\omega_I^{-1}\omega_J)^2=-1$。易错之处：不要以为
$\Im\tau$ 的标度在条件中起作用（它在 $\omega^{-1}F$ 中约掉）；也不要漏掉
「$F$ 与 $\hat N$ 同时反号」的第二个解 $F=-(\Im\tau)\omega_J$、
$\hat N=-I$，它对应把 Chan-Paton 线丛换成其逆（即相反定向的膜），正文
括号中一句带过。
\end{yizhu}

\begin{yizhu}[含糊点拨]
按：「没有瞬子修正」的论证值得慢慢看。边界条件 (11.6) 加上 $K$ 中的
瞬子方程，把法向导数 $\partial_1\Phi$ 用切向导数 $\partial_0\Phi$ 表示
出来，于是 (11.7) 变成 $(\omega-iF)\partial_0\Phi=0$。由于
$\omega-iF=-i\Omega_I\Im\tau$ 非退化（$\Omega_I$ 是全纯辛形式），只得
$\partial_0\Phi=0$；即瞬子在边界上为常数，再由 $K$-解析性（全纯映射的
唯一延拓）得它处处为常数。这里用的瞬子方程属于复结构 $K$，与 (11.8)
中作用量里的 $\omega$、$F$ 一致；不要误以为论证用的是 $I$ 或 $J$。
另外，「填满整个空间的膜」指 $Y=X$ 的情形，此时边界条件只在弦的端点处
施加，与膜内部的拓扑无关。
\end{yizhu}

\begin{yizhu}[含糊点拨]
按：为什么只作「部分层化」（只对 $T^*{\cal V}$、${\cal V}\subset\M$
型开集作层化）？因为完整的形变量子化层化在精确理论中会失效：
Moyal-Wigner 乘积是 $\hbar$ 的无穷级数，对一般的全纯函数并不收敛。原文
请读者试算 $1/(x-a)$ 与 $1/(y-b)$ 的 Moyal-Wigner 乘积来体会这一点：这两
个函数各自在 $\Bbb C^2$ 的适当开集上全纯，但乘积的无穷级数每一阶都给出
$1/(x-a)^m(y-b)^n$ 型项而不能截断，在极点附近没有意义。限制到
$T^*{\cal V}$ 并只取沿纤维多项式增长的函数后，代数由次数 $\le 1$ 的元素
生成；次数为 $m$ 与 $n$ 的元素之积若出现更低次项，须带 $\hbar$ 的正幂，
而 $\hbar$ 的次数为 $1$，故二阶以上的修正项无处安放。于是整个代数被二阶
数据完全确定，微扰论自动截断，不再有收敛问题。这正是原文说「一切都由二
阶发生的事情决定」的意思。
\end{yizhu}

\begin{yizhu}[物理图像]
按：c.c. 膜的 $(A,B,A)$ 性质可以这样记。在复结构 $K$ 中，$\omega_K$ 是
实的 $(1,1)$ 形式，而 $F=(\Im\tau)\omega_J$ 关于 $K$ 是
$(2,0)\oplus(0,2)$ 型的，故余各向同性条件 $(\omega^{-1}F)^2=-1$ 成立，
其中 $\hat N=\omega_K^{-1}\omega_J=I$ 是复结构。在复结构 $J$ 中，$\omega_J$
是 $(1,1)$ 形式，$F$ 因而是 $(1,1)$ 型，故规范场满足 $B$-膜（全纯）条件。
在复结构 $I$ 中，$\omega_I$ 是 $(1,1)$ 形式，$\omega-iF=-i\Omega_I\Im\tau$
是 $(2,0)$ 型的全纯辛形式，故它又是余各向同性 $A$-膜，条件即 (11.5)。从
二维场论看，c.c. 弦的代数就是 $\M$ 上的微分算子层
${\cal D}_{K_\M^{1/2}}$；它的用处是提供几何 Langlands 纲领所需的
${\cal D}$-模：$({\cal B}_{c.c.},{\cal B}_{c.c.})$ 弦给出环层，其它复结构
$K$ 中的 $A$-膜给出该环层上的模层，这些模层就是 $\M$ 上的 ${\cal D}$-模
（参见 Witten 2008 综述的相应讨论）。
\end{yizhu}

\begin{yizhu}[易错点]
按：时间反演把代数映到相反代数，这一步只用到「反转边界上算子插入的
顺序」，与四维的 $\EUT$ 如何作用在场上无关；后者只用来保证该对称性在
$t=1$ 的扭化理论中确实存在。由 ${\cal D}_{\cL}^{op}\cong
{\cal D}_{\cL^{-1}\otimes K_X}$，自相反的条件是
${\cal L}\cong{\cal L}^{-1}\otimes K_X$，即 ${\cal L}^2\cong K_X$，
故取 ${\cal L}=K_X^{1/2}$。注意这里「自相反」是代数层的同构，而不是
把 $A$-模型映到 $B$-模型的镜像对称；$K^{1/2}$ 的出现纯属代数上的
要求，与自旋结构的存在性无关（$\M$ 恰是自旋流形，故平方根存在）。
\end{yizhu}

\begin{yizhu}[原文疑误]
按：末段「the sheaf of c.c. strings is the sheaf of differential
operators on $\MH$」一语中的 $\MH$ 应为 $\M$。理由有三：其一，紧接着
说 $\hat w$ 是 $\M$ 上的三阶微分算子、其符号是 $\MH$ 上全纯函数的
$p_\alpha p_\beta p_\gamma$ 项；其二，前文刚刚确定的层是
${\cal D}_{K_{\EUM}^{1/2}}$，它是作用在 $\M$ 的典范丛平方根上的；
其三，本段标题的「整体代数」是取 $\M$ 的整体截面，而非 $\MH$ 的。
原文此处的 $\MH$ 应是笔误（或排版时漏掉了 $H$ 的下标改法），译文
照录，请读者按 $\M$ 理解。
\end{yizhu}

\begin{yizhu}[物理图像]
按：整体代数与量子可积系统的对应是本节的落点。$\MH$ 上 Hitchin
纤维化的对易哈密顿量，即 $\int_C\alpha\,{\cal P}(\varphi)$ 型的整体
全纯函数，经典上相互对易；把它们量子化为 $\M$ 上相互对易的微分算子，
就得到 $\M$ 上的量子可积系统。Beilinson--Drinfeld 的证明来自二维共形
场论在水平 $k=-h$ 的流代数（$h$ 为对偶 Coxeter 数，即 $k$ 取临界
水平），这一点与几何 Langlands 中「Hecke 特征层」的构造同源。原文
指出这些函数不能量子化为除 $K_\M^{1/2}$ 以外任何线丛上的微分算子，
说明这种量子化不是随意选取线丛就能做到的；而 $K_\M^{1/2}$ 正是时间
反演对称性挑出的那个。
\end{yizhu}
